\section{Module categories and Morita constraints}\label{sec:newalgebra}
These criteria use module categories forced by algebra objects and
complement the central-induction system. They do not impose a pivotal
structure on the original category. Field-automorphism and tensor-sign
reductions are kept with the arithmetic reductions in
Section~\ref{sec:reductions}.

\subsection{An involution-module obstruction}
\begin{theorem}[Antiinvariant module representation]\label{thm:involution}
Let $R$ be a commutative fusion ring and let $g\ne1$ be an invertible basis
element with $g^2=1$. Suppose the set $J=\{x:gx=x\}$ is nonempty. If $R$
has a complex fusion categorification, there are integer matrices
$T_i\in M_{|J|}(\ZZ)$, indexed by all basis elements, satisfying
\begin{align}
 T_0&=I,&T_iT_j&=\sum_kN_{ij}^kT_k,&T_{i^*}&=T_i^{\mathsf T},\label{eq:module-rep}\\
 |(T_i)_{xy}|&\le N_{ix}^{y},&
 (T_i)_{xy}&\equiv N_{ix}^{y}\pmod2 &&(x,y\in J).\label{eq:module-bound}
\end{align}
Failure of this finite system is a general categorification obstruction.
\end{theorem}
\begin{proof}
Let $\mathcal C$ be a hypothetical categorification. The additive category
generated by a fixed simple $x$ is a rank-one module category over the
pointed subcategory generated by $g$. Choose isomorphisms $g\otimes x\simeq x$.
The scalar module-associativity equation says that the pointed associator
three-cocycle is a coboundary. Thus that subcategory is tensor equivalent
to $\Vec_{C_2}$, and $A=1\oplus g$ admits the regular connected separable
algebra structure. This is a use of module associativity, not an assumption
that $g$ has a central lift. Standard separable-algebra module theory is
as in \cite[Chapter~7]{EGNO}.

Let $\mathcal M=\mathcal C_A$ be the category of right modules. Its free
and forgetful functors $I_A,F_A$ satisfy
$F_AI_A(X)=X\oplus(X\otimes g)$. Adjunction shows that a two-element
$g$-orbit gives one simple module, whose restriction is the sum of the two
simples in the orbit. A fixed simple $x$ gives two nonisomorphic simples
$M_x^+,M_x^-$, each restricting to $x$. Indeed,
$\dim\operatorname{End}(I_A(x))=2$, so the squared multiplicities in its
semisimple decomposition sum to two. The mixed Gram entries exclude
constituents from other orbits. Every simple module occurs in a free
module because $A$ is separable.

The algebra automorphism $\alpha$ acting as $+1$ on $1$ and $-1$ on $g$
induces an involution $Q$ of $\mathcal M$ by twisting the right action.
It fixes the module attached to a two-element orbit and interchanges
$M_x^+$ and $M_x^-$. The latter assertion follows directly because the
nonzero action map $x\otimes g\to x$ changes sign, whereas an
endomorphism of the simple underlying object is scalar. This twisting
functor commutes with the left action of $\mathcal C$.

Consequently the $(-1)$-eigenspace of $Q$ on $K_0(\mathcal M)\otimes\QQ$
is a representation of $R$, with basis
$e_x=[M_x^+]-[M_x^-]$, $x\in J$. Between two fixed pairs the action matrix
has a block $\left(\begin{smallmatrix}a&b\\b&a\end{smallmatrix}\right)$.
Forgetting modules gives $a+b=N_{ix}^{y}$, while its action on the
antiinvariant basis has coefficient $a-b$. Nonnegativity and integrality
of $a,b$ give \eqref{eq:module-bound}. The module associator gives the
multiplication identity in \eqref{eq:module-rep}. Finally module
Frobenius reciprocity makes the action of $i^*$ the transpose of the
action of $i$. All $e_x$ have the same squared length two, so the
restricted matrices obey the same transpose identity. No spherical or
unitary structure has been used.
\end{proof}

\begin{proposition}[Exact generator enumeration]\label{prop:module-enum}
The conditions of Theorem~\ref{thm:involution} can be decided by a finite
integer search. If a selected set of basis elements generates $R\otimes\QQ$
as an algebra, it suffices to enumerate their matrices, reconstruct the
other matrices by rational polynomial expressions, and check
\eqref{eq:module-rep}--\eqref{eq:module-bound}.
\end{proposition}
\begin{proof}
Every entry has an explicit finite interval and parity. Choose monomials
$w_0=1,w_1,\ldots,w_{r-1}$ forming a rational basis of the rank-$r$ algebra,
and let $V$ have as its rows their coefficients in the based-ring basis.
The expression for $b_i$ is row $i$ of $V^{-1}$ applied to the monomials.
Thus any representation is uniquely reconstructed from the selected
generator matrices. Testing all the ring relations, transpose identities,
and entry bounds after reconstruction accepts every possible solution
and nothing outside the stated necessary system.
\end{proof}

The direct enumeration implementation treats $|J|\le2$ and chooses one
or two generators when such a set is found. Applying it to the census
gives nine complete exclusion certificates, including $\RR{3}{79}$ and
$\RR{3}{216}$. The respective generator searches contain $4608$ and
$1536$ assignments. These are automatic applications of one theorem;
their indices are not inputs to the test. The spectral refinement in
Section~\ref{sec:spectral-involution} also treats larger fixed-point sets
and supplies two further exclusions with $|J|=4$.


\subsection{Spectral pruning of the antiinvariant module}
\label{sec:spectral-involution}
The matrix system in Theorem~\ref{thm:involution} is useful beyond the
two-dimensional antiinvariant space. Its entry bounds impose an additional
spectral restriction that makes larger fixed-point sets accessible.

\begin{proposition}[Perron-factor exclusion]\label{prop:module-perron}
Retain the notation of Theorem~\ref{thm:involution}, and let $i=i^*$.
Suppose the nonnegative regular fusion matrix $N_i$ is irreducible.
Then no $T_i$ can have an eigenvalue whose minimal polynomial contains
$\FPdim(i)$ as a root. Moreover, if for every $x\in J$ there are
$a,y\in J$ with $N_{ax}^{y}$ odd, then $T_g=I$.
\end{proposition}
\begin{proof}
Write $B_i=(N_{ix}^{y})_{x,y\in J}$. Since $J$ is a proper subset of
the basis, the strict Perron--Frobenius inequality for a proper principal
submatrix of an irreducible nonnegative matrix gives
\[
 \rho(T_i)\leq\rho(|T_i|)\leq\rho(B_i)<\rho(N_i)=\FPdim(i).
\]
The characteristic polynomial of the integer matrix $T_i$ is rational;
thus an irreducible factor containing $\FPdim(i)$ cannot divide it.
The bounds on $T_g$ make it diagonal with entries $\pm1$.
For $a\in J$, the identity $ga=a$ gives $T_gT_a=T_a$.
A diagonal entry $-1$ would therefore force the corresponding row of
every such $T_a$ to vanish, contradicting the indicated odd entry.
\end{proof}

If $i$ generates $R\otimes\mathbb Q$, express each basis element as a
rational polynomial $p_a(i)$. The preceding proposition removes the
Perron factor from the possible minimal polynomial of $T_i$; when
$T_g=I$, every remaining factor must also divide $p_g(t)-1$.
Since $T_i\equiv B_i\pmod2$, its characteristic polynomial must reduce
to that of $B_i$ modulo two. Enumerate all products of the remaining
irreducible factors of total degree $|J|$, allowing repetitions. Each
product fixes both $\operatorname{Tr}(T_i)$ and
$\operatorname{Tr}(T_i^2)=\sum_{x,y}(T_i)_{xy}^2$.
These equalities reduce the finite entry search substantially.
Reconstruction by the $p_a$ and verification of the entire system
\eqref{eq:module-rep}--\eqref{eq:module-bound} finish the test.
The search is exhaustive because every possible module representation
has one of the enumerated rational characteristic polynomials.

\begin{proposition}\label{prop:two-four-fixed-exclusions}
The rings $\RR{8}{2945}$ and $\RR{8}{3071}$ have no categorification
as complex fusion categories.
\end{proposition}
\begin{proof}
In both rings the involution is $g=b_1$ and
$J=\{b_2,b_3,b_4,b_5\}$; all basis elements are self-dual.
All tensor identities used below are checked automatically from the
census tensors in the supplementary certificates.

For $\RR{8}{2945}$, the element $b_3$ generates the rational algebra,
its regular matrix is irreducible, and
\[
 \det(tI-N_3)=t(t-2)(t^2-23t-1)(t^2+t-1).
\]
The Perron factor is $t^2-23t-1$. The parity criterion gives $T_g=I$,
and the character with $b_3$-value zero has $g$-value $-1$, so this
character cannot occur. Thus only $t-2$ and $t^2+t-1$ remain.
The fixed block has characteristic polynomial
$t^4+t^2+1$ modulo two. Of the three degree-four products of the two
remaining factors, only $(t^2+t-1)^2$ has this reduction. Hence
\[
 \operatorname{Tr}(T_3)=-2,\qquad
 \operatorname{Tr}(T_3^2)=6,\qquad T_3^2+T_3-I=0.
\]
In the ordered basis $J$, the required odd entries of $T_3$ are the
last two diagonal entries and the symmetric pairs $(1,4),(2,3)$.
Their squared contributions already sum to six. Consequently
\[
 T_3=\begin{pmatrix}
 0&0&0&\varepsilon\\
 0&0&\delta&0\\
 0&\delta&-1&0\\
 \varepsilon&0&0&-1
 \end{pmatrix},\qquad \varepsilon,\delta\in\{1,-1\}.
\]
Rational reconstruction, reduced modulo $t^2+t-1$, gives
$T_2=-3T_3-4I$. Its first diagonal entry is $-4$, contradicting the
bound $|(T_2)_{11}|\leq N_{b_2b_2}^{b_2}=2$.

For $\RR{8}{3071}$, the regular matrix $N_3$ is irreducible and
\[
 \det(tI-N_3)=t^3(t+2)(t^2-26t+24).
\]
The last factor is the Perron factor. After removing it, the symmetric
matrix $T_3$ has eigenvalues in $\{0,-2\}$, so $T_3^2=-2T_3$.
Every diagonal entry of its fixed block is six, and therefore every
diagonal entry $d$ of $T_3$ is even. The diagonal of the matrix identity
gives
\[
 d^2+\sum_{y\ne x}(T_3)_{xy}^2=-2d.
\]
It follows that $d\in[-2,0]\cap2\mathbb Z=\{-2,0\}$, and either
choice forces all off-diagonal entries in that row to vanish. This
contradicts $N_{b_3b_2}^{b_3}=7$, which requires the corresponding
off-diagonal entry to be odd.
\end{proof}

The exact generator search checks respectively four and zero
trace-and-norm-compatible matrices for these two records. The second
search uses $b_5$ as a rational algebra generator; the displayed direct
argument uses $b_3$ to make the contradiction shorter. Both certificates
are general categorification obstructions, with no pivotal assumption.


\subsection{Descent to a rank-three near-group subring}
\begin{theorem}[Rank-three subring obstruction]\label{thm:rank3-subring}
Suppose the Grothendieck ring of a complex fusion category contains
three distinct basis elements $1,g,x$ satisfying
\[
 g^2=1,\qquad gx=xg=x,\qquad x^2=1+g+mx.
\]
Then $m\leq2$.
\end{theorem}
\begin{proof}
The three basis elements are closed under products and duals, so their
simple objects generate a fusion subcategory with Grothendieck ring
$K(1,0,m,0)$. By \cite[Theorem~4.1]{Ostrik3}, categorifiability of this
ring implies $m\leq2$. The cited theorem uses the Galois-pseudounitary
reduction for this near-group family and does not impose a pivotality
hypothesis on the ambient category.
\end{proof}

In $\RR{4}{162}$, the basis elements $(b_0,b_3,b_4)$ have precisely
these products with $m=4$. The exact certificate checks all coefficients
in their nine products, including the vanishing coefficients outside
this subring. Theorem~\ref{thm:rank3-subring} therefore excludes this
ring from general complex categorification. The same automatic scan is
applied to the entire census through multiplicity eight.

\subsection{Small-dimension Morita obstructions}\label{sec:small-morita}
The next criteria use categorical Morita equivalence rather than additional
indicator-height searches. They apply to general complex fusion categories.
In particular, no chosen pivotal structure enters their hypotheses.

\subsubsection{A cyclic-index bound and dimension thirty}
For group-theoretical data $(G,\omega,H,\psi)$, the simple objects of
$\mathcal C(G,\omega,H,\psi)$ are indexed by a double coset $HgH$ and
an irreducible projective representation $\pi$ of
$S_g=H\cap gHg^{-1}$. Their Frobenius--Perron dimensions are
\begin{equation}\label{eq:gt-fp-new}
 [H:S_g]\dim\pi.
\end{equation}
The objects supported on $H$ form a fusion subcategory equivalent to
$\Rep(H)$. One can see this after replacing the data by an equivalent
adapted-cocycle presentation with $\psi=1$ and
$\omega|_{G\times G\times H}=1$; see
\cite[Proposition~2.5]{SchGT}. These descriptions also follow from
\cite{GNgt}.

\begin{lemma}[Cyclic-index bound]\label{lem:cyclic-index}
Suppose a finite group $G$ has a cyclic normal subgroup $B$ of index at
most two. Every simple object of every group-theoretical category
$\mathcal C(G,\omega,H,\psi)$ has Frobenius--Perron dimension at most two.
\end{lemma}
\begin{proof}
Put $K=H\cap B$. Since every subgroup of a cyclic group is characteristic,
$K$ is normal in $G$. Thus $K\subseteq S_g\subseteq H$ and
$[H:K]\leq2$. If $[H:S_g]=2$, then $S_g=K$ is cyclic. Every irreducible
projective representation of a cyclic group over $\CC$ has dimension one,
so \eqref{eq:gt-fp-new} is two.

If $S_g=H$, it suffices to bound the dimensions of projective
representations of $H$. In the twisted group algebra of $H$, the
subalgebra on $K$ is commutative and semisimple: its cyclic generator
satisfies $u^{|K|}=c\ne0$, a polynomial with distinct roots. The remaining
coset, when present, permutes its weight spaces in orbits of size at most
two. In an irreducible module there is one such orbit. For an orbit of
size one, a coset generator squares to a nonzero scalar and hence has a
one-dimensional invariant subspace. For an orbit of size two, a line in
one weight space and its translate span an invariant two-dimensional
subspace. Irreducibility therefore gives dimension at most two. This
argument includes arbitrary projective factor sets.
\end{proof}

\begin{proposition}[Integral dimension thirty]\label{prop:dimension30}
Every simple object of an integral fusion category of Frobenius--Perron
dimension $30$ has dimension one or two.
\end{proposition}
\begin{proof}
The category is group-theoretical by
\cite[Theorem~9.2]{ENOWGT}. Its presenting group has order thirty.
Every group of order thirty is solvable: the possible numbers of Sylow
five- and three-subgroups are respectively $1,6$ and $1,10$; if neither
is one, the corresponding nonidentity elements already number $24+20>29$.
A normal Sylow subgroup and the solvable quotient prove the claim.
Hall's theorem now supplies a subgroup of order fifteen. It is cyclic
by the Sylow theorems and normal because it has index two.
Lemma~\ref{lem:cyclic-index} applies.
\end{proof}
The adjective \emph{integral} is essential: the statement does not claim
the same bound for nonintegral categories of total dimension thirty.
The certificate checks a positive integral dimension character, not just
an integer value of the total dimension.

\subsubsection{A Morita dichotomy in dimension sixty}
\begin{lemma}\label{lem:sixty-dichotomy}
A fusion category of Frobenius--Perron dimension sixty is either solvable
or categorically Morita equivalent to $\Vec_{A_5}^{\omega}$ for some
three-cocycle $\omega$.
If it is integral and solvable, has exactly two invertible simples, and
satisfies $\mathcal C_{\mathrm{ad}}=\mathcal C$, every simple has dimension
at most four.
\end{lemma}
\begin{proof}
By \cite[Theorem~9.16]{ENOWGT}, the category is weakly group-theoretical.
It is therefore Morita equivalent to a nilpotent category of the same
dimension, obtained from vector spaces by successive extensions by
finite groups whose orders multiply to sixty. If all these groups are
solvable, refining their normal series into cyclic prime-order factors
shows that the category is solvable. Solvability is invariant under
categorical Morita equivalence; see \cite[Section~4.2]{ENOWGT}.
Otherwise an extension group has nonsolvable order dividing sixty.
All proper divisors involve at most two primes, apart from thirty,
which was handled above. Burnside's theorem and the fact that the unique
nonabelian simple group of order sixty is $A_5$ imply that this group is
$A_5$ itself. Its order uses the entire dimension. The nilpotent category
is consequently pointed, of the form $\Vec_{A_5}^{\omega}$.

For the second assertion, a nontrivial last step of a solvable
construction is a cyclic prime-order extension or equivariantization.
An extension would give a nontrivial grading, contrary to
$\mathcal C_{\mathrm{ad}}=\mathcal C$. In an equivariantization
$\mathcal C=\mathcal D^{C_p}$, the canonical $\Rep(C_p)$ is a fusion
subcategory. Exactly two invertible simples force $p=2$, and
$\FPdim(\mathcal D)=30$.

The simples of a $C_2$-equivariantization are indexed by orbits of simples
of $\mathcal D$ and projective representations of their stabilizers.
Those stabilizers are trivial or cyclic of order two, so the projective
representations have dimension one. A simple dimension in $\mathcal C$
is therefore either $d$ or $2d$, with $d$ a simple dimension in
$\mathcal D$. Since $\mathcal C$ is integral, $d$ is rational; as an
algebraic integer it is an integer. Thus $\mathcal D$ is integral.
Proposition~\ref{prop:dimension30} gives $d\leq2$, proving the bound four.
\end{proof}

\subsubsection{Cocycle-independent involution returns}
For a normalized three-cocycle on a group of order two, its value on the
triple of nonidentity elements is unchanged by a normalized coboundary.
In particular, triviality of its restriction class implies that this
value is one.

\begin{theorem}[Involution-return indicator]\label{thm:involution-returns}
Let $\mathcal C(G,\omega,H,\psi)$ be group-theoretical, let $g\in G$,
and put $S=H\cap gHg^{-1}$. Suppose every element in
\[
 T_g=\{r\in gH:r^2\in S\}
\]
is an involution conjugate in $G$ to an element of $H$.
For the simple object $M(g,\pi)$ associated to an irreducible projective
representation $\pi$ of $S$, one has
\begin{equation}\label{eq:involution-return}
 \nu_2(M(g,\pi))=\frac{|T_g|\dim\pi}{|S|}.
\end{equation}
In particular, if $T_g\ne\varnothing$, these simples are self-dual.
This conclusion holds for every admissible $\omega,\psi$.
\end{theorem}
\begin{proof}
Use the equivalent adapted presentation supplied by
\cite[Proposition~2.5]{SchGT}; its cocycle restricts to one on $H^3$.
If $r$ is an involution conjugate into $H$, the restriction class on
$\langle r\rangle$ is trivial. The preceding coboundary observation gives
$\omega(r,r,r)=1$ in this adapted presentation, independently of any
remaining values of the cocycle.

The second-indicator specialization of
\cite[Theorem~4.2, equation~(4.6)]{SchGT} is
\[
 \nu_2(M(g,\pi))=
 \frac1{|S|}\sum_{r\in T_g}\pi_{-2}(r)\chi_\pi(r^{-2}).
\]
For an involution, the scalar $\pi_{-2}(r)$ in that formula is
$\omega(r,r,r)^{-1}=1$, whereas
$\chi_\pi(r^{-2})=\chi_\pi(1)=\dim\pi$.
This proves \eqref{eq:involution-return}. Group-theoretical categories
have their positive spherical structures; a simple with nonzero second
indicator is self-dual. Self-duality is preserved under tensor
equivalence, so the conclusion is independent of the presentation.
\end{proof}

\begin{corollary}[Dihedral supports in $A_5$]\label{cor:a5-returns}
In every category $\mathcal C(A_5,\omega,D_{2q},\psi)$ with
$q\in\{3,5\}$, every simple of dimension $q$ has second indicator one.
In particular all such simples are self-dual.
\end{corollary}
\begin{proof}
The dihedral subgroup has order $2q$. Its ordinary irreducible
representations have dimensions one or two, so a simple of dimension
$q>2$ is not supported on it. Formula~\eqref{eq:gt-fp-new} then forces
stabilizer order two and a one-dimensional projective representation.
For these double cosets the following exact group data hold:
\[
\begin{array}{c|c|c|c}
 |H|&\text{double-coset sizes}&|S|&|T_g|\text{ when }|S|=2\\\hline
 6&6,18,36&2&2\\
 10&10,50&2&2
\end{array}
\]
Every element in $T_g$ is an involution. All fifteen involutions in $A_5$
are conjugate, and $H$ contains an involution. Thus
Theorem~\ref{thm:involution-returns} gives $\nu_2=2/2=1$.

These group assertions are certified independently from the sixty even
permutations of five letters. The verifier generates every order-$2q$
subgroup using an element of order $q$ and an involution, checks that they
form one conjugacy class, and enumerates every right coset and every
square-return element. It does not test only a selected coset or an
assumed cocycle. The coset lists are included in each certificate.
\end{proof}

\begin{theorem}[Tensor-checkable small-dimension obstruction]\label{thm:small-morita}
Let $R$ be an integral fusion ring. Each of the following conditions
excludes a general complex categorification.
\begin{enumerate}[label=\textup{(\alph*)}]
\item $\FPdim(R)=30$ and some basis dimension exceeds two.
\item $\FPdim(R)=60$, the adjoint based subring is $R$, and exactly two
basis elements are invertible. Some basis dimension exceeds four.
For a prime $q\in\{3,5\}$ the complete set of based-subring dimensions is
$\{1,2,2q,60\}$, and a basis element of dimension $q$ is not self-dual.
\end{enumerate}
\end{theorem}
\begin{proof}
Part (a) is Proposition~\ref{prop:dimension30}.
For (b), Lemma~\ref{lem:sixty-dichotomy} excludes solvability and gives a
presentation $\mathcal C(A_5,\omega,H,\psi)$.
The canonical $\Rep(H)$ subcategory forces
$|H|\in\{1,2,2q,60\}$. If $|H|\leq2$, every simple dimension is at most
two, contradicting the stated dimension exceeding four. If $H=A_5$,
the category is $\Rep(A_5)$ and has only one invertible simple, because
$A_5$ is perfect. Thus $H$ is an order-$2q$ subgroup, necessarily a
conjugate of $D_{2q}$. Corollary~\ref{cor:a5-returns} contradicts the
non-self-dual dimension-$q$ basis element.
\end{proof}

\subsubsection{Uniform census application}
The implementation supplies the same tensor test at every multiplicity.
It first checks a positive integral dimension character; it then
computes the adjoint closure, all based subrings, duality, and the number
of invertibles. Illustrative applications are
\[
\begin{array}{c|c|c|c}
 \text{ring}&\text{sorted dimension type}&\text{branch}&q\\\hline
 \RR{2}{51}&(1,1,2,2,5,5)&\text{(b)}&5\\
 \RR{3}{15}&(1,1,2,3,3,6)&\text{(b)}&3\\
 \RR{4}{19}&(1,1,1,1,1,5)&\text{(a)}&\text{--}
\end{array}
\]
The program finds these applications from the tensor, without selecting
labels in advance. The current census scan and its exact dimension
checks are recorded by \filename{software/morita_controls.py}.
The criterion also excludes $\RR{2}{25}$ and $\RR{2}{29}$.
Overlap with other criteria is counted only once in the sequential
exclusion table.

\subsection{Two-sided module Gram constraints}\label{sec:dual-module-gram}
Let $\mathcal C$ be a fusion category, let $A$ be a connected separable algebra,
and let $\mathcal M=\mathcal C_A$. Write $M_i$ for the action matrix of a
simple $X_i\in\mathcal C$, and $\mathcal D=\operatorname{Fun}_{\mathcal C}
(\mathcal M,\mathcal M)$ for the dual category, with a fixed choice of
opposite tensor convention. All dimension statements here use
Frobenius--Perron dimensions, not a positive pivotal structure.

\begin{proposition}[Two-sided module Gram identities]\label{prop:dual-module-gram}
For every simple $m\in\mathcal M$, let $B_m$ be the matrix of the evaluation
functor $\mathcal D\to\mathcal M$, $Y\mapsto Y(m)$. Then
\begin{equation}\label{eq:module-gram}
 B_mB_m^{\mathsf T}=\sum_i(M_i)_{mm}M_i.
\end{equation}
The tensor unit column of $B_m$ is $e_m$. If $d_n$ are the positive module
dimensions in any fixed normalization, each column $v$ belonging to a
simple $Y\in\mathcal D$ satisfies
\begin{equation}\label{eq:module-dim}
 \sum_n v_nd_n=d_m\FPdim(Y).
\end{equation}
\end{proposition}
\begin{proof}
The internal endomorphism algebra of $m$ in $\mathcal C$ is
$A_m=\bigoplus_i(M_i)_{mm}X_i$, up to the harmless replacement of every
index by its dual; these multiplicities agree by module reciprocity.
The module category is equivalently $\mathcal C_{A_m}$, sending $m$ to
the free rank-one module. The dual category is the category of
$A_m$-bimodules. Evaluation becomes the functor forgetting the left action,
whose adjoint is free left multiplication by $A_m$. Their composite on
$\mathcal M$ is therefore the action of $A_m$. Adjunction gives
\eqref{eq:module-gram}. Evaluation of the tensor unit is $m$, giving the
specified column. The commuting action of the dual category on the
positive module dimension vector has eigenvalue $\FPdim(Y)$; this gives
\eqref{eq:module-dim}. These assertions are standard consequences of
separable-algebra Morita theory; see \cite[Chapter~7]{EGNO}.
\end{proof}

\subsubsection{The split-involution input and a finite system}
Suppose $g^2=1$ is an invertible object fixing a simple object of
$\mathcal C$. The module-associativity argument in the main manuscript
supplies the algebra $A=1\oplus g$. Free two-element $g$-orbits give one
simple of $\mathcal M$; fixed simples give two. The sign automorphism of
$A$ interchanges the two actions on every fixed simple. Thus the action
matrices $M_i$ can be reconstructed from the full tensor and an
anti-invariant representation satisfying the parity, transpose and entry
bounds of Theorem~\ref{thm:involution}.

Normalize the dimension of the free module $A$ to one. An orbit module
has dimension $\FPdim(X)$ and a split fixed module has dimension
$\FPdim(X)/2$. These module dimensions need not be algebraic integers.
At $m=A$, the Gram matrix in \eqref{eq:module-gram} is $I+M_g$.
Since $M_g$ is an involutive permutation, its nonnegative integral Gram
columns are forced: a fixed point has two unit columns, and a transposed
pair has one column supported on that pair. Hence
\eqref{eq:module-dim} determines the multiset of simple dimensions of
$\mathcal D$. One of the unit-dimension columns is its tensor unit.

For another simple $m$, remove its known tensor-unit column. For every
remaining prescribed dual dimension $e$, enumerate nonzero integer
vectors $v\ge0$ with
\[
 v_i^2\le Q_{ii},\quad v_iv_j\le Q_{ij},\quad
 \sum_i v_i d_i=e d_m,
 \qquad Q=\sum_i(M_i)_{mm}M_i-e_me_m^{\mathsf T}.
\]
The integer multiplicities of these columns must add their outer
products to $Q$ and respect the prescribed number of columns of each
dual dimension. Relaxing the multiplicities to nonnegative rationals
gives a necessary rational linear system. An exact Farkas vector proves
infeasibility. A general exclusion requires infeasibility for every
admissible anti-invariant representation; selecting only one is invalid.

The executable implementation treats at most two fixed simples and
simple-dimension fields of degree at most four. It enumerates all
admissible generator matrices and reconstructs the complete representation.
The vector search uses exact field coordinates and explicit integer
bounds. A domain-budget limit or an ineligible fixed-point count gives
no conclusion. Negative checking, when a certificate is proposed,
reconstructs the full domain and checks only exact arithmetic.


The implementation in \filename{software/module_dual_gram.py} is a
uniform necessary-condition test. It does not assign categorical
existence when the relaxed system is feasible. Its successful
certificates, if any, must enter the same exact-checking and disjoint
allocation process as every other obstruction.

\subsection{Coupling all Morita action matrices}\label{sec:full-morita-gram}
The two-sided module test can be strengthened by retaining the complete
associativity pairing, rather than its restrictions to one module simple at
a time. The strengthening uses Frobenius--Perron positivity only to recognize
module categories; it makes no pivotal or spherical assumption.

\begin{theorem}[Full Morita action Gram identity]\label{thm:full-morita-gram}
Let ${}_{\mathcal C}\mathcal M_{\mathcal D}$ be an invertible bimodule
category, with simple objects $m_0,\ldots,m_{s-1}$. Write
\[
 L_i(a,b)=[X_i\otimes m_a:m_b],\qquad
 R_j(a,b)=[m_a\otimes Y_j:m_b]
\]
for its left and right action matrices. Then
\begin{equation}\label{eq:full-morita-gram}
 \sum_j R_j(a,d)R_j(b,c)
       =\sum_i L_i(b,a)L_i(c,d)
       \qquad(0\leq a,b,c,d<s).
\end{equation}
Consequently the matrix
\[
 G_{(a,d),(b,c)}=\sum_i L_i(b,a)L_i(c,d)
\]
has a nonnegative integral Gram factorization whose columns are the
vectorized matrices $R_j$. Each such column belongs to
$\operatorname{im}_{\mathbb Q}G$, and its coordinate indexed by $(a,d)$ is
at most $\lfloor\sqrt{G_{(a,d),(a,d)}}\rfloor$.
The corresponding matrices commute with every $L_i$, include the identity,
and are permuted by transposition.
\end{theorem}
\begin{proof}
Use duals in the Morita bicategory. The multiplicity of $Y_j$ in
$m_b^*m_c$ is $R_j(b,c)$, by Frobenius reciprocity. Thus the multiplicity
of $m_d$ in $m_a(m_b^*m_c)$ is the left-hand side of
\eqref{eq:full-morita-gram}. The multiplicity of $X_i$ in $m_am_b^*$ is
$L_i(b,a)$, so associativity identifies the same multiplicity with its
right-hand side. If $W$ is the matrix of vectorized right actions, this
is $G=WW^{\mathsf T}$. Hence
$\operatorname{im}_{\mathbb R}W=\operatorname{im}_{\mathbb R}G$; since all
matrices are rational, each column of $W$ belongs to the stated rational
image. A diagonal entry of $G$ is the sum of the squares of the
corresponding entries of $W$, which proves the bound. The commuting
left and right actions, the tensor unit of $\mathcal D$, and
$R_{j^*}=R_j^{\mathsf T}$ prove the remaining assertions.
\end{proof}

Suppose now that $A=\mathbf1\oplus g$ is the split involution algebra and
$m_0=A$ in $\mathcal C_A$. The existence argument for this algebra and the
exhaustive module-action construction are those of the preceding module
criterion. Since the internal end of $m_0$ is $A$, setting $c=d=0$ in
\eqref{eq:full-morita-gram} gives
\begin{equation}\label{eq:morita-first-columns}
 \sum_j R_j(\,\cdot\,,0)R_j(\,\cdot\,,0)^{\mathsf T}=I+L_g.
\end{equation}
Here $L_g$ is an involutive permutation matrix. Its right-hand side has
a unique nonnegative integral Gram factorization with nonzero columns,
up to permutation of columns: a fixed point contributes two copies of its unit vector, and
a transposition contributes the sum of its two unit vectors. Indeed,
zero off-diagonal entries prevent columns from meeting different
orbits. On a fixed point the squared entries sum to two; on a
transposition the diagonal entries and the off-diagonal entry are all
one. These observations prove the assertion orbit by orbit.
Transposition supplies the same finite list for the possible first rows
of each right action matrix.

The implementation intersects the rational image of $G$ with these
prescribed first rows and columns. Choosing free coordinates among the
actual matrix entries, the diagonal bounds in
Theorem~\ref{thm:full-morita-gram} give a finite integer box. Every
candidate is checked for nonnegativity, the entrywise necessary bounds
$vv^{\mathsf T}\leq G$, and commutation with the left actions. If the
candidates are $v_1,\ldots,v_t$, the unknown multiplicities must satisfy
\[
 \sum_{h=1}^t z_hv_hv_h^{\mathsf T}
       =G-\operatorname{vec}(I)\operatorname{vec}(I)^{\mathsf T},
 \qquad z_h\in\mathbb Z_{\geq0},
\]
together with the first-column multiplicities from
\eqref{eq:morita-first-columns}, after removing the unit action.
An exact Farkas certificate for the relaxation $z_h\geq0$ excludes that
module-action branch. As before, exclusion of the fusion ring requires
exhausting every eligible branch. The program
\filename{software/full_morita_gram.py} reconstructs the finite domains
and checks separating vectors using integer arithmetic.

\begin{proposition}[The involution dual of the first residual]
\label{prop:first-residual-morita-self-ring}
If $\mathcal C$ categorifies $\RR{3}{215}$ and $g$ is its nontrivial
invertible simple, then $\mathbf1\oplus g$ admits a connected separable
algebra structure and
\[
 K_0\bigl({}_A\mathcal C_A\bigr)\cong K_0(\mathcal C)
 \quad\text{as based rings},\qquad A=\mathbf1\oplus g.
\]
This assertion does not identify the two fusion categories up to tensor
equivalence.
\end{proposition}
\begin{proof}
The involution fixes two simples, so the previously established split
involution argument applies. Exhaustive reconstruction gives sixteen
module-action branches. Two violate the rank bound
$\operatorname{rank}G\leq\operatorname{rank}K_0({}_A\mathcal C_A)$,
and twelve have exact Farkas certificates for the full Gram system.
Each of the remaining two has Gram rank six and six dual simples.
Therefore its six right action matrices are linearly independent; in
particular they cannot repeat. Exhaustive selection from the finite
candidate lists gives exactly one Gram frame for each branch, up to
permuting its columns. Multiplying those six matrices reconstructs a
nonnegative integral fusion tensor, and a recorded based permutation
identifies all its coefficients with $\RR{3}{215}$.

The complete sixteen-branch record is
\filename{research/new_frontier/frontier_full_morita.json}.
The script
\filename{research/new_frontier/verify_frontier_full_morita.py}
reconstructs the module actions and candidate domains, checks all rank
and Farkas obstructions, proves uniqueness of each surviving frame by
finite exact enumeration, and verifies every matrix product and the
based-ring identification. Verification does not call an optimizer.
\end{proof}

The two surviving frames are based bimodule data, not associators or
categorifications. In particular the proposition does not supply an
algebra $\mathbf1\oplus X$ for a noninvertible simple $X$.
The uniform full Gram scan covers all $47$ residual inputs through
multiplicity eight: $21$ are eligible for the implemented involution
search and $26$ are ineligible. It yields no additional whole-ring
exclusion. All $54$ supplied positive models pass the control scan.
These outcomes are retained separately from the exclusion ledger.

\subsection{Integral multiplicities in the Morita Gram identity}
\label{sec:integral-morita-rotation}

The full Morita action Gram identity of
Theorem~\ref{thm:full-morita-gram} has an elementary integral refinement.
The unknown attached to a possible right-action matrix counts simple
objects of the dual category; it is an integer, even when distinct simple
objects have the same action matrix. Retaining this integrality detects
obstructions invisible in the rational nonnegative relaxation.

\begin{proposition}[Morita Gram congruence obstruction]
\label{prop:morita-gram-congruence}
Fix one of the exhaustively enumerated module actions for a forced split
involution algebra $\mathbf1\oplus g$. Let $Ax=b$ be its full Morita Gram
system, including the prescribed first-column multiplicities, with all
possible right-action matrices retained. If there exist $q>1$ and an
integral row vector $y$ such that
\[
 yA\equiv0\pmod q,
 \qquad yb\not\equiv0\pmod q,
\]
then this module action cannot occur in a complex fusion category.
\end{proposition}
\begin{proof}
Every categorification would give a vector $x$ of nonnegative integral
multiplicities satisfying $Ax=b$. Multiplication by $y$ would give
$yb=yAx\equiv0\pmod q$, a contradiction. All entries of $A$ and $b$ come
from the integral fusion and module actions. The argument therefore
requires no pivotal or spherical structure and no positivity assumption
on categorical dimensions.
\end{proof}

For discovery, Gaussian elimination over $\mathbb F_2$ produces the row
vector $y$. Verification only multiplies integer vectors and checks the
two congruences. Repeated equality rows may be removed by a deterministic
rule. The completeness of the right-action candidate list is exactly the
finite Gram enumeration already used in
Section~\ref{sec:full-morita-gram}; no new search cutoff is introduced.

\begin{proposition}\label{prop:round3-higher-stabilizer}
The rings \texttt{R6\_5\_638} and \texttt{R6\_7\_1225} admit no
categorification by a fusion category over $\mathbb C$.
\end{proposition}
\begin{proof}
The unrestricted signed reduction is verified separately for each full
fusion tensor. There are respectively $11\,800$ and $46\,856$ signed
tables, partitioned into $418$ and $1\,560$ sign orbits. Each has one
original gauge orbit. Every other orbit is excluded by its exact
character obstructions; within the original orbit the only admissible
nonzero dimension characters lie in the Frobenius--Perron Galois orbit.
It therefore suffices to exclude pseudo-unitary categorification.

In each case the involution-module enumeration has sixteen branches.
Two fail the full Morita rank bound. Twelve fail
Proposition~\ref{prop:morita-gram-congruence}, with $q=2$. Both remaining
branches force negative stabilizer commutators on the two fixed simple
objects, indexed $4$ and $5$ in the stored tensor. Their cubic rotation
spaces have common dimension $n=4$ for \texttt{R6\_5\_638} and $n=6$
for \texttt{R6\_7\_1225}. The higher-dimensional Clifford rotation
criterion of Section~\ref{sec:stabilizer-rotation} applies. The exponent
prime supports are respectively $\{2,5,13\}$ and $\{2,5,29\}$, so the
third indicators belong respectively to
\[
 \{-2,1\}\qquad\text{and}\qquad\{-3,0\}.
\]
The actual separating certificates already hold on the weaker intervals
$[-2,1]$ and $[-3,3/2]$; they do not rely on selecting a discrete indicator
value.

For either tensor the complete low-induction enumeration gives $91$
frames, of which $31$ have a nonempty character-labelling set. The $50$
labellings and four second-indicator choices yield $200$ exact
height-six systems. Every system has a verified integral separating
functional. The two dimension fields are $\mathbb Q(\sqrt{65})$ and
$\mathbb Q(\sqrt{145})$; their low-column counts are respectively
$(217,225,225,61,61)$ and $(469,481,481,127,127)$. All possible twist
orders are covered by the inherited relative-trace construction.
Thus the two remaining module branches also fail.
\end{proof}

The exact verifier is
\filename{software/round3_integral_morita.py}, with certificate kind
\texttt{general\_integral\_morita\_rotation}. The full finite domains,
congruence vectors, separating vectors, and optimizer-free replay reports
are stored in \filename{research/round3/higher/}.

\subsection{An order-six Morita--grading obstruction}
\label{sec:morita-order6}

The dimension reconstruction above can force a dual category whose
based ring has a short exhaustive parametrization. This supplies two
further exclusions without assuming a pivotal structure on the original
category.

\begin{theorem}\label{thm:morita-order6}
Let $\mathcal C$ be a fusion category, and let $g$ be an invertible
involution fixing a simple object. Set $A=\one\oplus g$. Suppose the
simple Frobenius--Perron dimensions of
$\mathcal D={}_A\mathcal C_A$ are
\[
  \underbrace{1,\ldots,1}_{6},\qquad
  \underbrace{d,\ldots,d}_{6},\qquad d=2+\sqrt5.
\]
Then $\mathcal C$ has a faithful $C_3$ grading. In particular, a fusion
ring with trivial universal grading and the stated forced dual dimension
profile cannot be categorified. This excludes
$\RR{5}{100}$ and $\RR{5}{408}$ over $\CC$.
\end{theorem}

\begin{proof}
The algebra $A$ exists by the fixed-object argument of
Theorem~\ref{thm:involution}. We first exhaust the possible based rings
of $\mathcal D$. Its invertible objects form a group $G$ of order six.
The multiplicity of an invertible $h$ in $\rho\rho^*$, for a
noninvertible simple $\rho$, is one precisely when $h\rho\cong\rho$.
The identity $d^2=1+4d$ and irrationality of $d$ therefore imply that
the stabilizer of $\rho$ is trivial. Thus the noninvertibles are exactly
$h\rho$, $h\in G$. The right action is free for the same reason, so
\begin{equation}\label{eq:morita-six-relations}
 \rho h=\phi(h)\rho,\qquad
 \rho^2=u+\sum_{h\in G}a_hh\rho,\qquad
 \sum_{h\in G}a_h=4,
\end{equation}
where $\phi\in\operatorname{Aut}(G)$, $u\in G$, and
$a_h\in\ZZ_{\ge0}$. Associativity gives
$\phi^2=\operatorname{Ad}(u)$ and $\phi(u)=u$; comparison of
$(\rho^2)h$ with $\rho(\rho h)$ also makes the coefficient function
invariant under
\[
 k\longmapsto\phi^2(h)k\phi(h)^{-1}.
\]
Only $C_6$ and $S_3$ can occur as $G$.

\smallskip\noindent
\emph{The cyclic case.}
The automorphism $\phi$ is either the identity or inversion. Inversion
would make $(a_h)$ constant on cosets of the square subgroup of order
three, contradicting $\sum a_h=4$. Hence $\phi$ is the identity.
Use additive labels $0,\ldots,5$ for $C_6$. Duality gives
$\rho^*=(-u)\rho$ and
\begin{equation}\label{eq:morita-six-domain}
 u\in\ZZ/6\ZZ,\qquad a_k\ge0,\qquad
 \sum_{k=0}^5a_k=4,\qquad a_k=a_{u-k}.
\end{equation}
There are exactly $60$ pairs $(u,a)$ satisfying these conditions.
Each supplies the entire candidate tensor by
\[
 h\cdot k=h+k,\qquad h\cdot(k\rho)=(h+k)\rho,\qquad
 (h\rho)(k\rho)=h+k+u+\sum_j a_j(h+k+j)\rho.
\]
An exact exhaustive calculation partitions the $60$ pairs as follows:
\begin{center}
\begin{tabular}{@{}rl@{}}
\toprule
Number & Exact consequence \\
\midrule
$42$ & A formal-codegree norm obstruction \\
$6$ & A closed based subring isomorphic to $\RR{2}{20}$ \\
$12$ & A faithful $C_3$ grading with neutral invertible involution \\
\bottomrule
\end{tabular}
\end{center}
The first row is a general obstruction here: by
\cite[Lemma~3.1]{SchTwo}, a category with these two distinct
Frobenius--Perron dimensions is Galois conjugate to a pseudo-unitary
category with the same based ring. In that category $\FPdim(\mathcal
D)/f$ must be an algebraic integer for every formal codegree $f$.
The cyclic candidate ring is commutative, so the Casimir characteristic
roots used in the check are its formal codegrees, with no representation
dimension correction. Its global dimension is
\[
 D=6+6d^2=60+24\sqrt5,
 \qquad p_D(t)=t^2-120t+720.
\]
For each of the $42$ cases, the certificate gives the exact minimal
polynomial $q$ of a formal codegree and a prime $p$ for which
\begin{equation}\label{eq:morita-normalized-norm}
 \frac{v_p(|p_D(0)|)}{\deg p_D}
 <\frac{v_p(|q(0)|)}{\deg q}.
\end{equation}
This contradicts integrality of $D/f$: take the norm in any common
number field containing $D$ and $f$. The second row uses the previously
established general exclusion for $\RR{2}{20}$, whose original exact
certificate is replayed, not just its allocated status.
For each of the remaining $12$ pairs the grading is explicitly
\begin{equation}\label{eq:morita-six-grading}
 \deg(h)=h\pmod3,\qquad \deg(h\rho)=h+2u\pmod3.
\end{equation}
Every nonzero structure constant respects this degree assignment.
The unique nonidentity invertible involution, namely $3\in C_6$, has
degree zero.

\smallskip\noindent
\emph{The nonabelian case.}
Every automorphism of $S_3$ is inner. Replacing $\rho$ by an invertible
multiple makes $\phi$ the identity. Then $u$ is central, so $u=1$,
and $a$ is constant on conjugacy classes. Its values
$(a_e,a_t,a_c)$ on the identity, transpositions and three-cycles satisfy
\[
 a_e+3a_t+2a_c=4.
\]
The complete domain and its obstructions are
\begin{center}
\begin{tabular}{@{}ll@{}}
\toprule
$(a_e,a_t,a_c)$ & Obstruction \\
\midrule
$(4,0,0)$ & The rank-two subring $\rho^2=1+4\rho$ \\
$(2,0,1)$ & The closed subring $\RR{2}{20}$ \\
$(0,0,2)$ & The closed subring $\RR{2}{21}$ \\
$(1,1,0)$ & A one-dimensional formal-codegree norm obstruction \\
\bottomrule
\end{tabular}
\end{center}
The first row is excluded by \cite{Ostrik2}. In the next two rows the
subring is generated by $C_3$ and $\rho$; its based isomorphism and the
previously established general certificate are checked exactly.
For the last row, extend the sign representation of $S_3$ by assigning
$\rho$ a root $\lambda$ of $\lambda^2+2\lambda-1=0$. Equation
\eqref{eq:morita-six-relations} verifies this as a one-dimensional
representation of the full ring. Its genuine formal codegree is
\[
 f=\sum_{h\in S_3}1+\sum_{h\in S_3}\lambda^2
   =6(1+\lambda^2)=24\pm12\sqrt2,
 \qquad q(t)=t^2-48t+288.
\]
Here $\rho^*=\rho$. Since $v_2(720)=4$ and $v_2(288)=5$, condition
\eqref{eq:morita-normalized-norm} holds. The two-dimension reduction
again makes this a general categorical obstruction. In particular, this
step does not identify an unweighted Casimir eigenvalue from a
higher-dimensional irreducible representation with a formal codegree.
All nonabelian possibilities have now been excluded.

\smallskip\noindent
\emph{Returning through Morita duality.}
Consequently $\mathcal D$ has one of the faithful gradings
\eqref{eq:morita-six-grading}. Regard $A$ as the distinguished simple
right module in $\mathcal M=\mathcal C_A$, and let
$B=\underline{\operatorname{End}}_{\mathcal D}(A)$. Double Morita
duality identifies $\mathcal C$ with ${}_B\mathcal D_B$, up to the
opposite tensor convention; see \cite[Chapter~7]{EGNO}.
For a simple $A$-bimodule $Y$, its multiplicity in $B$ is
\[
 \dim\operatorname{Hom}_{\mathcal M}(A\otimes_A Y,A)
 =\dim\operatorname{Hom}_{\mathcal M}(Y,A).
\]
There are exactly two simple bimodules restricting to the right module
$A$: the regular bimodule and its twist by the nontrivial character of
$C_2$. Each has multiplicity one. This can also be read directly from
the two unit columns at $A$ in the Gram matrix $I+M_g$ of
Proposition~\ref{prop:dual-module-gram}; off-diagonal entries in that
row vanish. All other simple bimodules have zero multiplicity. Thus
$B=\one\oplus b$, where $b$ is a nontrivial invertible involution.
In the cyclic case $b$ is necessarily $3\in C_6$, hence $B$ lies in
the neutral component.

A connected separable algebra in the neutral component of a faithful
group grading transfers that grading to its bimodule category. Indeed,
the left and right algebra actions preserve each homogeneous support,
and tensor product over the algebra multiplies the degrees. Every
component is nonzero: for a nonzero homogeneous object $Y$, the free
bimodule $B\otimes Y\otimes B$ is nonzero and has the same degree.
It follows that $\mathcal C$ is faithfully $C_3$-graded.

Finally, for $\RR{5}{100}$ the involution is basis label $4$, the
fixed labels are $3,5$, and the exhaustive anti-invariant generator
search has $16$ tuples. For $\RR{5}{408}$ these data are respectively
$1$, $4,5$, and $8$. Each search has exactly one admissible
representation, in which the involution acts as $+I$ on the two
anti-invariant coordinates. The Gram reconstruction forces exactly the
dimension profile in the statement. In both original tensors the
multiplicative closure of the supports of all $XX^*$ is the whole
basis. Thus the adjoint subcategory would be all of $\mathcal C$,
which rules out any nontrivial faithful grading. This proves the two
claimed exclusions.
\end{proof}

\paragraph{Exact finite certificate.}
The implementation is \filename{software/new_morita_order6.py}, with
audit output in
\filename{research/model_continuation/morita_order6_certificates.json}.
The input certificate stores the involution, fixed labels, generator
search size, every admissible anti-invariant representation, and the
reconstructed dual dimensions in exact field coordinates. Verification
reconstructs the exhaustive representation domain, rather than trusting
a proposed list of branches. The finite dual routine enumerates every
composition of $4$ into six nonnegative parts, applies exactly
\eqref{eq:morita-six-domain}, and checks that there are $60$ cases.
It constructs each full tensor, factors its Casimir polynomial exactly,
and checks the recorded norm witness, the full grading identities, or
the closed-subring support and unit-preserving based isomorphism.
The four $S_3$ class triples are independently exhausted. The inherited
certificates for $\RR{2}{20}$ and $\RR{2}{21}$ are replayed before they
are used. Thus the audit output is a reproducible record of a complete
finite necessary-condition calculation, with no unresolved branch
silently discarded.

A separate replay, \filename{research/model_continuation/replay_morita_order6.py},
checks equality with the saved finite audit and rejects mutations of the
involution, branch list, generator count, representation matrices, dual
dimension counts, and input basis labels. Its checks use explicit
exceptions and remain active under \texttt{python -O}.


\subsection{Forced pointed associators and the multiplicity-four frontier}
\label{sec:c4-forced-associator}

The next obstruction uses the associator of a pointed subcategory to restrict
central twists. It is a tensor criterion and applies to an entire family.

\begin{proposition}\label{prop:c4-forced-associator}
Let $m\geq0$, and suppose a fusion category has invertible objects
$1,t,t^2,t^3$, with $t$ of order four, and two self-dual noninvertible
objects $X,Y$ satisfying
\[
 \begin{gathered}
 tX=Y=Xt,\qquad tY=X=Yt,\\
 X^2=Y^2=1+t^2+m(X+Y),\qquad XY=t+t^3+m(X+Y).
 \end{gathered}
\]
Then the pointed associator has class $2$ in
$H^3(C_4,\mathbb C^\times)\cong\mathbb Z/4\mathbb Z$.
In a pseudo-unitary realization, every simple central summand of
$I(t)$ or $I(t^3)$ has twist satisfying $\theta^4=-1$.
\end{proposition}
\begin{proof}
Put $H=\langle t^2\rangle$. Because $t^2$ fixes a simple object,
the associator restricts trivially to $H$: the associativity scalars
of the action on that object give a two-cochain trivializing the
restriction. Thus its class in $H^3(C_4,\mathbb C^\times)$ is either
$0$ or $2$.

Suppose first that the class is zero. The regular pointed algebra
$A=1\oplus t\oplus t^2\oplus t^3$ exists. Its right-module category
has three simple objects. The invertible orbit contributes $A$;
the orbit $\{X,Y\}$ contributes two simples $M_+,M_-$, each with
underlying object $X\oplus Y$. Indeed its stabilizer is $C_2$,
whose projective irreducible representations are the two ordinary
one-dimensional representations, because $H^2(C_2,\mathbb C^\times)=0$.
This is the usual orbit--stabilizer description of modules over a
regular pointed algebra (equivalently, the classification of modules
over its action groupoid algebra).

Let $d=\operatorname{FPdim}(X)=\operatorname{FPdim}(Y)$, so
$d^2=2+2md$ and $d=m+\sqrt{m^2+2}$. Normalize module dimensions
by $\operatorname{FPdim}(A)=4$; their vector is $(1,d/2,d/2)$.
Every invertible fixes the simple module $A$ and permutes $M_+,M_-$.
On their difference line its action is a character
$\chi:C_4\to\{\pm1\}$, hence $\chi(t^2)=1$.
The action matrix of $X$ is symmetric, since $X$ is self-dual.
Its first row and column are $(0,1,1)$ by free--forgetful adjunction.
The dimension eigenvector therefore forces its remaining block to be
\[
 \begin{pmatrix}a&2m-a\\2m-a&a\end{pmatrix}
 \quad (a\in\mathbb Z).
\]
Thus its eigenvalue on the difference line is the integer
$\lambda=2a-2m$. Applying the fusion relation for $X^2$ on this
line gives
\[
 \lambda^2=2+m(1+\chi(t))\lambda.
\]
For $\chi(t)=-1$ this says $\lambda^2=2$; for $\chi(t)=1$
it says $(\lambda-m)^2=m^2+2$. Both are impossible for integers:
$m^2+2$ lies strictly between consecutive squares for $m\geq1$,
and equals $2$ for $m=0$. Hence class $0$ is impossible.

For completeness use representatives $a,b,c\in\{0,1,2,3\}$ and
standard cyclic cocycles
\[
 \omega_k(a,b,c)=\exp\left(\frac{2\pi i k}{4}
          a\left\lfloor\frac{b+c}{4}\right\rfloor\right).
\]
Their restrictions have $\omega_k(2,2,2)=(-1)^k$, confirming
that the restriction kernel is $\{0,2\}$. The canonical
pseudo-unitary indicator formula gives
$\nu_4(t)=i^k$ and $\nu_4(t^3)=i^{9k}$; for $k=2$ both equal
$-1$. Reversing the cocycle convention inverts these values and
has no effect on this conclusion.

The center indicator formula reads
\[
 \operatorname{FPdim}(\mathcal C)\nu_4(t)
   =\sum_{Z\subset I(t)}[Z:I(t)]\operatorname{FPdim}(Z)\theta_Z^4.
\]
The positive coefficients on the right sum to
$\operatorname{FPdim}(I(t))=\operatorname{FPdim}(\mathcal C)$.
Equality in the triangle inequality forces every $\theta_Z^4=-1$.
The same argument applies to $t^3$; compare
\cite[Proposition~3.8]{SchTwo}.
\end{proof}

\begin{lemma}\label{lem:c4-uniform-module}
Under the hypotheses of Proposition~\ref{prop:c4-forced-associator},
$m$ is even. For the algebra $1\oplus t^2$, its anti-invariant
module action on the two fixed-object coordinates is, up to sign gauge,
\[
 T_t=\begin{pmatrix}0&-1\\1&0\end{pmatrix},\qquad
 T_{t^2}=-I,\qquad T_X=T_Y=0.
\]
In particular both noninvertibles have the anticommuting stabilizer
used in Proposition~\ref{prop:stabilizer-rotation}.
\end{lemma}
\begin{proof}
The involution-module construction gives
$T_t=\left(\begin{smallmatrix}0&\epsilon\\\delta&0\end{smallmatrix}\right)$,
where $\epsilon,\delta\in\{\pm1\}$, and
$T_{t^2}=\epsilon\delta I$. Since $X,Y$ are self-dual,
$T_X$ and $T_Y$ are symmetric integer matrices; they commute with
$T_t$ and satisfy $T_Y=T_tT_X$.
If $\epsilon\delta=1$, commutation makes the diagonal entries
of $T_X$ equal. Its eigenvalue $\lambda$ on the $+1$ eigenspace
of $T_t$ is therefore an integer. But the relation for $X^2$
gives $\lambda^2=2+2m\lambda$, which is impossible as above.
Consequently $\epsilon\delta=-1$. A symmetric matrix commuting
with this skew-symmetric $T_t$ is scalar, so $T_X=aI$.
The symmetry of $T_Y=aT_t$ forces $a=0$. Finally every entry
of $T_X$ is congruent modulo two to the corresponding coefficient
of multiplication by $X$ on $\{X,Y\}$, namely $m$. Hence $m$
is even. Changing the sign of one fixed-object coordinate gives
the displayed sign gauge.
\end{proof}

\begin{theorem}\label{thm:c4-multiplicity-four}
The fusion ring $\RR{4}{6}$ has no categorification over $\mathbb C$.
\end{theorem}
\begin{proof}
Its full tensor has the family in
Proposition~\ref{prop:c4-forced-associator} with $m=4$ and
$(t,t^2,t^3,X,Y)=(b_1,b_5,b_2,b_3,b_4)$.
All simple Frobenius--Perron dimensions belong to
$\{1,4+3\sqrt2\}$. The two-dimension Galois-pseudo-unitarity
theorem therefore reduces general existence to pseudo-unitary
existence, without a pivotality assumption on the original category.

Lemma~\ref{lem:c4-uniform-module} gives exactly one anti-invariant
module representation for the split involution algebra $1\oplus t^2$,
after sign gauge. The finite module checker independently reconstructs it.
The action of $t^2$ is $-I$ on both fixed-object coordinates.
The cubic spaces of $X$ and $Y$ each have dimension four;
Proposition~\ref{prop:stabilizer-rotation} therefore forces
\[
 -2\leq\operatorname{Re}\nu_3(X),
          \operatorname{Re}\nu_3(Y)\leq1.
\]
The exact central-induction calculation has ten low frames and
four second-indicator assignments per frame. All low summands
have twist one. Every high restriction column containing $t$ or
$t^3$ is restricted by Proposition~\ref{prop:c4-forced-associator}
to $\theta^4=-1$. In the normalized relative-trace domain over
$\mathbb Q(\sqrt2)$, this means its fourth coordinate is exactly
$(-1,0)$. Keeping only those trace types is necessary; actual roots
with $\theta^4=-1$ have precisely that coordinate under every
relative Galois embedding. The order-four trace system includes
all root orders with the allowed prime support $\{2,3\}$,
including arbitrarily large orders through their zero trace type.

Each of the resulting forty rational nonnegative systems has a
recorded integral separating vector. Verification reconstructs
the low frames, the unique module branch, the entire joint trace
domain, the four second-indicator assignments, and the forced-twist
variable deletion. For every system $Az=b$, its vector $y$
satisfies $A^Ty\geq0$ and $b^Ty<0$. Therefore no feasible central
induction data, and hence no pseudo-unitary categorification,
exists. This contradicts the general-to-pseudo-unitary reduction.
\end{proof}

\paragraph{Certificate and implementation.}
The tensor-based implementation is
\filename{software/round3_smallest.py}. Its verifier never calls
an optimizer; the discovery routine merely proposes exact rays.
The certificate, proof audit, all-census applicability scan,
positive controls and optimized replay are retained in
\filename{research/round3/smallest/}.


\subsection{A stabilizer obstruction to scalar pivotal rotation}
\label{sec:stabilizer-rotation}

\begin{proposition}[Anticommuting stabilizers and third rotation]
\label{prop:stabilizer-rotation}
Let $\mathcal C$ be a spherical fusion category, let $g$ be an invertible
of order two for which $A=\mathbf1\oplus g$ is an algebra, and let $Y$
be a simple object with $gY\cong Y\cong Yg$. Suppose that left
multiplication by $g$ exchanges the two right $A$-module structures
on $Y$. Put $m=\dim\operatorname{Hom}(\mathbf1,Y^{\otimes3})$.
Then $m$ is even. If $m>0$, the pivotal third-rotation operator on this
space is not scalar. More precisely,
\[
 \nu_3(Y)=-(a+b\zeta_3+c\zeta_3^2),
 \qquad a,b,c\in\ZZ_{\ge0},\quad a+b+c=m/2.
\]
If $3$ is coprime to the Frobenius--Schur exponent, this becomes
\[
 \nu_3(Y)\in
 \{-m/2+3b:0\le b\le\lfloor m/4\rfloor\}.
\]
In particular, if $m=2$ under this coprimality hypothesis, then
\[
 \nu_3(Y)=-1.
\]
\end{proposition}
\begin{proof}
Use strict monoidal notation and an associative identification
$q:g^2\xrightarrow\sim\mathbf1$ supplied by the algebra $A$.
Choose isomorphisms $l:gY\to Y$ and $r:Yg\to Y$, normalized by
\[
 l(\operatorname{id}_g\otimes l)=q\otimes\operatorname{id}_Y,
 \qquad r(r\otimes\operatorname{id}_g)=\operatorname{id}_Y\otimes q.
\]
The two normalized right module structures are $r$ and $-r$.
Transporting the right action on $gY$ through $l$ gives
\[
 r'=l(\operatorname{id}_g\otimes r)(l^{-1}\otimes\operatorname{id}_g).
\]
The hypothesis says $r'=-r$, or equivalently
\begin{equation}\label{eq:stabilizer-anticommutator}
 r(l\otimes\operatorname{id}_g)=-l(\operatorname{id}_g\otimes r).
\end{equation}
This also shows directly how the hypothesis is read from the
anti-invariant module matrices: it is $T_g[Y,Y]=-1$.

Define an invertible endomorphism of $Y^2$ by
\[
 a=(r\otimes l)(\operatorname{id}_Y\otimes q^{-1}\otimes\operatorname{id}_Y),
\]
and put $P=a\otimes\operatorname{id}_Y$, $Q=\operatorname{id}_Y\otimes a$ on $Y^3$.
Expanding either composite inserts two copies of $q^{-1}$, with common
intermediate tensor word $YggYggY$. The endpoint maps agree. On the
middle $gYg$, the two composites use the two sides of
\eqref{eq:stabilizer-anticommutator}. Consequently
\[
 PQ=-QP.
\]
The restrictions to $V=\operatorname{Hom}(\mathbf1,Y^3)$ remain
invertible. Taking determinants shows that $\dim V$ is even.

Let $E$ be pivotal cyclic rotation moving the last $Y$ to the first
position. The canonical cyclic map
\[
 \operatorname{Hom}(\mathbf1,B\otimes Y)
 \longrightarrow \operatorname{Hom}(\mathbf1,Y\otimes B)
\]
is natural in the entire object $B$. Applying this naturality to
$B=Y^2$ and $a\in\operatorname{End}(B)$ gives the exact identity
\[
 EP=QE.
\]
The construction and its naturality are those of
\cite[Definition~3.1, Lemma~3.4 and Theorem~5.1]{NSpiv}.
No braiding or self-duality sign is used. If $E$ were scalar on nonzero
$V$, then $P=Q$, contradicting $PQ=-QP$ and invertibility.

The module-action normalizations also give $a^2=I$: after two pair
insertions, the first endpoint applies $r(r\otimes\operatorname{id}_g)$,
the second applies $l(\operatorname{id}_g\otimes l)$, and the two
resulting copies of $q$ cancel the inserted copies of $q^{-1}$ by
associativity. Set
$A_j=E^jPE^{-j}$ for $j\in\ZZ/3\ZZ$. These are involutions, and the
anticommutation of $A_0=P$ and $A_1=Q$, together with cyclic conjugation,
makes them pairwise anticommuting. They therefore give a representation
of the complex Clifford algebra $\mathrm{Cl}_3\cong M_2(\CC)\oplus
M_2(\CC)$. In that algebra the element
\[
 U=\frac{-I+A_0A_1+A_1A_2+A_2A_0}{2}
\]
satisfies $U^3=I$, $UA_jU^{-1}=A_{j+1}$, and has trace $-1$ on each
two-dimensional simple module. These identities also follow by direct
multiplication using $A_i^2=I$ and $A_iA_j=-A_jA_i$.
The eight-dimensional universal algebra calculation is reproduced with
rational coefficients by
\filename{research/round2_audit/clifford_rotation.py}; it verifies
$U^3=I$ and $UA_j=A_{j+1}U$. Its scalar coefficient is $-1/2$, while
the three bivectors are traceless on either two-dimensional simple,
giving trace $-1$ without a matrix-choice assumption.
Thus $S=U^{-1}E$ commutes with the Clifford action and satisfies
$S^3=I$. On its multiplicity spaces, of total dimension $m/2$, let
$a,b,c$ be the multiplicities of $1,\zeta_3,\zeta_3^2$ for $S$.
Taking the trace of $E=US$ proves the asserted formula. Reversing the
cyclic convention exchanges $b$ and $c$, so the same value set applies
to the standard indicator convention.

On the other hand, the central-twist
formula puts $\nu_3(Y)$ in $\QQ(\zeta_N)$, where $N$ is the
Frobenius--Schur exponent. If $3\nmid N$, then
$\QQ(\zeta_N)\cap\QQ(\zeta_3)=\QQ$. Hence $b=c$, and
$a+2b=m/2$ gives the rational-value list.
\end{proof}

\begin{remark}[No self-duality requirement]
The stabilizer-rotation proposition applies equally to a non-self-dual
simple $Y$. Its hypotheses and proof use isomorphisms $gY\cong Y\cong Yg$,
not an isomorphism $Y\cong Y^*$. In particular, rigidity gives
\[
 m=\dim\operatorname{Hom}(\mathbf1,Y^{\otimes3})=N_{Y,Y}^{Y^*};
\]
one must not replace the last index by $Y$ unless $Y$ is self-dual.
The pair-insertion operators act on $Y^3$ itself, and naturality of pivotal
rotation applies to $Y^2$ and its endomorphism $a$ without dualizing $Y$.
Thus the same Clifford multiplicity formula applies to each member of a
dual pair separately.
\end{remark}

\begin{corollary}[Projected third-indicator spectrum]
\label{cor:stabilizer-average}
Under the hypotheses of Proposition~\ref{prop:stabilizer-rotation},
\[
 -m/2\le\operatorname{Re}\nu_3(Y)\le m/4,
 \qquad
 \operatorname{Re}\nu_3(Y)\in
 \{m/4-3a/2: a=0,1,\ldots,m/2\}.
\]
Let $K$ be any number field satisfying
$K\cap\mathbb Q(\zeta_3)=\mathbb Q$. The normalized relative trace of
$\nu_3(Y)$ to $K$ equals its real part and therefore belongs to the
displayed finite set. This includes every real number field and the
complex cyclotomic field $\mathbb Q(\zeta_8)$.
If instead $K$ contains $\zeta_3$, the unprojected possibilities are the
full set
\[
 \{-(a+b\zeta_3+c\zeta_3^2):
       a,b,c\in\mathbb Z_{\ge0},\ a+b+c=m/2\}.
\]
\end{corollary}
\begin{proof}
The preceding proposition gives
$\nu_3(Y)=-(a+b\zeta_3+c\zeta_3^2)$ with $a+b+c=m/2$.
Its real part is $-a+(b+c)/2=m/4-3a/2$, which gives the stated set and
bounds. Under the intersection hypothesis, $K(\zeta_3)/K$ is quadratic
and its nontrivial automorphism exchanges $\zeta_3$ and $\zeta_3^2$.
Consequently its normalized relative trace sends both roots to $-1/2$,
and sends the indicator to $-a+(b+c)/2$. Transitivity gives the same
value when the trace is taken from any larger common field.
The quadratic field $\mathbb Q(\zeta_3)=\mathbb Q(\sqrt{-3})$ does not
embed into $\mathbb Q(\zeta_8)$: the latter's three quadratic subfields
are $\mathbb Q(i)$, $\mathbb Q(\sqrt2)$ and
$\mathbb Q(\sqrt{-2})$. This proves the stated complex-field example.
When $K$ contains $\zeta_3$, normalized trace fixes the entire indicator,
so the original Clifford multiplicity set must be retained.
\end{proof}


\begin{corollary}\label{cor:frontier-third-indicators}
Every pseudo-unitary categorification of $\RR{3}{215}$, in the basis
of dimensions $(1,1,d,d,d-1,d+1)$, must satisfy
\[
 \nu_3(b_4)=\nu_3(b_5)=-1.
\]
\end{corollary}
\begin{proof}
The exact full Morita-action calculation for $A=1+g$ leaves two branches,
both with $T_g=-I$ on the two fixed simples $b_4,b_5$.
They satisfy $N_{44}^{4}=N_{55}^{5}=2$, and both are self-dual.
Their third-rotation spaces therefore have dimension two.
The global dimension is $136+32\sqrt{17}$, whose norm is
$1088=2^6\cdot17$. The Cauchy theorem
\cite[Theorem~3.9]{BNRW} excludes the prime $3$ from the
Frobenius--Schur exponent. Apply
Proposition~\ref{prop:stabilizer-rotation}.
\end{proof}

\subsection{Exact certificate and exhaustive coverage}
\label{subsec:stabilizer-rotation-certificate}
The new obstruction is implemented without a distinguished census label in
\filename{software/stabilizer_rotation.py}. Its input is the multiplication
tensor.  The certificate first reconstructs the complete list of
anti-invariant actions for the forced algebra $\mathbf1\oplus g$.
For the multiplicity-three frontier this list has sixteen entries.  Twelve
are rejected by exact separating functionals for the full Morita-action Gram
system, and two by its rank bound.  In each of the remaining two entries,
left multiplication by $g$ exchanges both pairs of module structures on the
fixed simples $b_4,b_5$.  Equivalently, the diagonal entries of $T_g$ on their
anti-invariant lines are both $-1$.  Their cubic invariant spaces have
 dimension two.  The pair-insertion theorem therefore forces
$\nu_3(b_4)=\nu_3(b_5)=-1$, since the Cauchy prime support is $\{2,17\}$.

The central calculation is exhaustive; it does not assume the rank-twenty-two
modular datum retained elsewhere as a relaxation witness.  All ninety-one
low induction frames are reconstructed from the tensor and its formal
codegrees.  Thirty-one frames admit character labellings; there are fifty
labellings in total and four possible second-indicator assignments.  The
certificate contains an exact Farkas vector for each of these two hundred
systems.  Each system includes the induction Gram equations, relative twist
traces through order three, cyclic-rotation character constraints, joint
Galois compatibility of dimensions and eigenvalues, and the low-row modular
relations.  It additionally imposes the two forced third indicators.  If
$Ax=b$ is such a system, its recorded integral vector $y$ is checked to satisfy
\[
 A^{\mathsf T}y\geq0,\qquad b^{\mathsf T}y<0.
\]
Every feasible central decomposition supplies $x\geq0$, so these inequalities
exclude every branch.  The variables allow all roots of unity whose prime
support is $\{2,17\}$ via their finite joint relative-trace types; no finite
upper bound on the actual exponent or central rank is imposed.  For clarity, the order-three relative-trace domain has exactly ten joint
patterns.  Write a root's order as $2^a17^b$.  If $b\geq2$ or $a\geq3$,
its normalized traces at powers $1,2,3$ all vanish.  Otherwise
$a\in\{0,1,2\}$ and $b\in\{0,1\}$ give the remaining nine distinct
patterns, with the two quadratic signs when $17$ divides the order.
For primitive seventeenth roots the first trace is
$(-1\pm\sqrt{17})/16$; the quadratic character of $2$ is $+1$ and
that of $3$ is $-1$.  This elementary enumeration explicitly includes
arbitrarily large exponent orders through the all-zero pattern.
Integral high-column decompositions are not needed for this final exclusion.

The general-scope wrapper is
\filename{software/galois_pseudo_reduction.py}. It reconstructs the full
unrestricted signed-sphericalization enumeration, including all $1480$
signed quotient tables and all sign orbits.  Every nongauge character
component is rejected by its exact arithmetic certificate, and the only
nonzero dimension-character orbit of the original ring is checked directly
to be its Frobenius--Perron orbit.  Consequently every hypothetical complex
categorification has a pseudo-unitary Galois conjugate to which the preceding
obstruction applies.  No positivity normalization is imposed during the
signed enumeration.  The combined certificate is stored in
\filename{certificates/stabilizer_rotation_215.json}; the public replay entry
point is \texttt{stabilizer\_rotation.verify\_general(R,c)}.

\begin{theorem}[Resolution of the multiplicity-three frontier]
\label{thm:frontier-noncategorifiable}\label{thm:stabilizer-rotation-frontier}
The based ring $\RR{3}{215}$ has no complex fusion categorification.
\end{theorem}
\begin{proof}
The unrestricted signed-sphericalization calculation of
Proposition~\ref{prop:frontier-galois-pu}, independently reconstructed by
the general certificate above, shows that any hypothetical
categorification has a pseudo-unitary Galois conjugate. Work with that
conjugate.

The complete involution-module domain has sixteen branches. The full
Morita-action Gram identities exclude fourteen, and either remaining
branch satisfies the hypotheses of
Corollary~\ref{cor:frontier-third-indicators}. Consequently both
$\nu_3(b_4)$ and $\nu_3(b_5)$ equal $-1$.

The exhaustive central-induction calculation in
Section~\ref{subsec:stabilizer-rotation-certificate} excludes those two
equalities for every permitted low frame, character assignment and
second-indicator choice. Each of its $200$ integral separating vectors
is checked against the reconstructed rational system; no numerical
infeasibility output is used as a certificate. Thus no pseudo-unitary
categorification exists, contradicting the reduction.
\end{proof}

This proof uses no identification of a hypothetical Drinfeld center with
a known braided category. The compatible cyclic-eight modular datum from
the previous continuation remains a useful test of the older necessary
conditions, but the local stabilizer and rotation constraints prevent it,
and every other central-induction branch, from arising here.

\subsection{Further stabilizer exclusions at multiplicity four}
\label{subsec:stabilizer-rotation-416}
\begin{corollary}\label{cor:stabilizer-416}
The ring $\RR{4}{416}$ has no complex fusion categorification.
\end{corollary}
\begin{proof}
The same local theorem excludes \texttt{R6\_4\_416} without changing its
hypotheses to fit the input.  The unique relevant invertible involution is
$b_1$, with fixed simples $b_4,b_5$.  Both exhaustive anti-invariant module
actions have negative stabilizer sign on both fixed objects.  Each cubic
invariant space has dimension four.  The Clifford-rotation trace therefore
lies in $\{-2,1\}$, since the Cauchy prime support is $\{2,5\}$.
In particular its upper bound is one.

The quadratic central-induction calculation reconstructs all $895$ low
frames.  After character reciprocity and joint Galois pruning, two relevant
character labellings remain.  All sixteen second-indicator choices for each
labelling are excluded by exact Farkas rays, using the interval bound
$-2\leq\nu_3(b_4),\nu_3(b_5)\leq1$.  Thus only $32$ separating vectors are
needed, and no discrete guess for the third indicators is required.

For general scope, the unrestricted signed enumeration has $2376$ quotient
tables in $83$ sign orbits.  All nongauge branches are rejected by $222$
exact character obstructions.  Of the original nonzero dimension-character
components, the non-Frobenius--Perron component has global dimension polynomial
$x^2-32x+128$ and fails the norm condition at prime $5$, required by the
Frobenius--Perron global dimension polynomial $x^2-800x+3200$.
Consequently the same general-to-Galois-pseudo-unitary reduction applies.
The combined general-scope certificate is
\texttt{certificates/stabilizer\_rotation\_416.json}; it has been replayed
under optimized Python with all numerical optimizer entry points disabled.

\end{proof}

\paragraph{The ring $\RR{4}{174}$.}
The four invertible basis elements have Frobenius--Perron dimension $1$;
the remaining two have dimension $d=4+3\sqrt2$, with $d^2=2+8d$.
Thus the global dimension is $72+48\sqrt2$, whose norm is $576$.
The involution-module calculation has a single antiinvariant branch, and
its stabilizer commutator is negative on both noninvertible simples.
Each cubic invariant space has dimension four. Since the global-dimension
prime support contains $3$, the exact relative traces of their third
indicators are allowed to be $-2,-1/2,$ or $1$. The computation first uses
the weaker interval $[-2,1]$.

At trace height six, there are ten low induction frames. Eight character
labellings survive the earlier exact conditions; their full second-indicator
assignments give 128 systems. Every system admits an exact separating ray,
already with the interval constraint. Hence no pseudo-unitary
categorification exists. The two distinct Frobenius--Perron dimensions
supply the previously established general-to-Galois-pseudo-unitary reduction,
so the exclusion holds for arbitrary complex fusion categories.
The composite certificate uses the general higher-stabilizer verifier,
with the reduction payload \texttt{two\_fp\_values}; all matrix and branch
data are reconstructed from the input tensor.

\paragraph{The ring $\RR{4}{427}$.}
Write $d=7+5\sqrt2$, so that $d^2=1+14d$. The Frobenius--Perron
dimensions are $(1,1,d,d,d-1,d+1)$ and the global dimension is
$400+280\sqrt2$, of norm $3200$. There is one antiinvariant
involution-module branch. Its stabilizer commutator is negative on the
last two simples, whose cubic invariant spaces both have dimension four.
The prime $3$ does not divide the global-dimension norm; the Clifford
criterion therefore restricts both third indicators to $\{-2,1\}$.

The trace-height-six computation enumerates 315 low induction frames.
The surviving character assignments comprise eight labellings and 32
second-indicator systems. Exact separating rays exclude every system
even using the weaker interval $[-2,1]$. This proves pseudo-unitary
nonexistence. The independent unrestricted signed-quotient reduction
has 2608 tables, forming 87 sign-gauge orbits, with 239 nongauge
character rejections. Its exact character certificates reduce every
categorification to the Galois-pseudo-unitary case. Combining the two
verified components gives general nonexistence through the same
\texttt{general\_stabilizer\_rotation\_higher} verifier.

\subsection{Ambient cyclotomic traces for quartic dimension fields}
\label{sec:cyclotomic-stabilizer}
The relative-trace method does not require the dimension field to be
quadratic. One may instead embed it into a cyclotomic field and retain
all coefficients of the resulting normalized traces.

\begin{proposition}[A finite ambient trace domain]
\label{prop:ambient-trace-domain}
Let $L=\QQ(\zeta_F)$ contain all simple dimensions of a pseudo-unitary
fusion category, and let $S$ be the prime support supplied by the
categorical Cauchy theorem. Suppose every prime dividing $F$ lies in
$S$. For height $h$, put
\[
 B=\prod_{p\in S}p^{\max\{v_p(F),1\}+\lfloor\log_p h\rfloor}.
\]
Every joint normalized relative-trace vector
\[
 \left(\frac{\operatorname{Tr}_{L(\theta)/L}(\theta^j)}
 {[L(\theta):L]}\right)_{j=1}^{h}
\]
of a root of unity whose order has prime support in $S$ is either zero,
or occurs among roots of unity of order dividing $B$.
\end{proposition}
\begin{proof}
If $p^a$ divides the order with
$a>\max\{v_p(F),1\}+\lfloor\log_p h\rfloor$, then for each
$1\le j\le h$ the $p$-primary part of $\theta^j$ has order strictly
larger than $p^{\max\{v_p(F),1\}}$. Summing over the last degree-$p$
cyclotomic layer gives zero: its conjugates differ by the complete set
of $p$th roots of unity. This layer fixes $L$ and the other primary
components, so transitivity of trace gives zero in $L$. Otherwise the
order divides $B$. Since $F\mid B$, all remaining traces can be computed
inside $\QQ(\zeta_B)$ by averaging the automorphisms
$\zeta_B\mapsto\zeta_B^a$ with $a\equiv1\pmod F$.
\end{proof}

Use each pair consisting of an admissible high restriction column $v$
and a trace vector as a nonnegative count variable. The Gram equations
remain integral. Apply normalized relative trace to the central
indicator formula of \cite[Theorem~4.1]{NS}; the low induction objects
have twist one. Expanding in the power basis of $L$ gives rational
linear equations for the first three moments. The usual second
indicator choices are exhausted. When $3\nmid F$ and $3\notin S$, the
third indicator lies both in $\QQ(\zeta_3)$ and in a cyclotomic field
of conductor coprime to three, hence is rational. Its cyclic
multiplicities give the necessary description
\[
 \nu_3(Y)=a-b/2,\qquad a,b\ge0,\qquad
 a+b=\dim\operatorname{Hom}(\mathbf1,Y^3).
\]
The count $b$ includes both primitive third-root eigenspaces.
For a four-dimensional cubic space with the anticommuting stabilizer
of Proposition~\ref{prop:stabilizer-rotation}, impose in addition
$\nu_3(Y)\le1$. These conditions are linear after adding nonnegative
slack variables. Allowing rational multiplicities enlarges the
necessary system and is therefore legitimate for exclusion.

\begin{theorem}[Quartic stabilizer-trace exclusions]
\label{thm:quartic-stabilizer}
The rings $\RR{4}{258}$, $\RR{4}{325}$, $\RR{8}{570}$ and
$\RR{8}{2203}$ have no complex
categorification.
\end{theorem}
\begin{proof}
The tensor-only anti-invariant module enumeration has one possible
representation for each ring. The order-two invertible has mixed
stabilizer sign $-1$ on the self-dual simple $b_5$ in the two
multiplicity-four rings, and on $b_4$ in the two multiplicity-eight
rings. Each of these cubic invariant spaces has dimension four.
Thus the preceding bound applies in every
pseudo-unitary categorification.

For $\RR{4}{258}$, set $u=2\cos(2\pi/20)$. The exact dimension-field
embedding gives
\[
 D=12u^3+20u^2-20u-20,\qquad
 p_D(t)=t^4-120t^3+2400t^2-16000t+32000.
\]
The Cauchy prime support is $\{2,5\}$. At height three,
Proposition~\ref{prop:ambient-trace-domain} gives $B=40$ and $31$
distinct trace vectors, including zero. The complete low-induction
search has three frames and four second-indicator choices per frame.
Every one of these twelve systems has an exact separating ray.

For $\RR{4}{325}$, set $u=2\cos(2\pi/16)$. Here
\[
 D=8u^3+16u^2,\qquad
 p_D(t)=t^4-128t^3+2560t^2-16384t+32768.
\]
The prime support is $\{2\}$, the height-three order bound is $32$,
and there are $25$ trace vectors. The two low frames and four
second-indicator choices per frame give eight exact separating rays.
For both $\RR{8}{570}$ and $\RR{8}{2203}$, take
$u=2\cos(2\pi/20)$. The shared dimension-field data are
\[
 D=160u^3+300u^2-220u-400,\qquad
 p_D(t)=t^4-1400t^3+44000t^2-400000t+800000.
\]
Again $S=\{2,5\}$, $B=40$, and the trace domain has $31$ vectors.
Each ring has seven low frames. There is one second-indicator choice
for $\RR{8}{570}$ and four for $\RR{8}{2203}$, giving respectively
seven and twenty-eight exact separating rays.
These computations establish pseudo-unitary nonexistence in all four
cases.

The general certificates also reconstruct all signed sphericalization
quotients and their character obstructions. For $\RR{4}{258}$ there
are $32$ quotients in one sign-gauge orbit. For $\RR{4}{325}$ there are
$96$ quotients in four orbits; the three nongauge orbits require nine
character rejections. The two multiplicity-eight rings give
respectively $96$ quotients in six orbits with ten nongauge character
rejections, and $1472$ quotients in $46$ orbits with $111$ such
rejections. The surviving orbit in each case is gauge
isomorphic to the original ring, and the nonpositive original
characters are excluded by the recorded exact tests. Consequently any
complex categorification would have a Galois conjugate with the
positive spherical dimensions just excluded.

The canonical implementation is
\filename{software/cyclotomic_stabilizer_trace.py}. Its
\texttt{verify\_general} entry point checks both the sphericalization
reduction and the complete pseudo-unitary certificate. It reconstructs
the module domain, field embedding, trace domain, low frames, high
columns and second-indicator alternatives from the input tensor; no
census identifier enters a mathematical test. Every retained high
column satisfies the exact integral dimension quotient and central
simple $d$-number conditions. All separating inequalities are checked
with integers. The certificate files are
\filename{certificates/cyclotomic_stabilizer_R6_4_258.json},
\filename{certificates/cyclotomic_stabilizer_R6_4_325.json},
\filename{certificates/cyclotomic_stabilizer_R6_8_570.json} and
\filename{certificates/cyclotomic_stabilizer_R6_8_2203.json}.
\end{proof}

The same ambient-field computation applies to the other quartic inputs
recorded in \filename{research/round2_literature/quartic/}. Its exploratory
modular refinement keeps global Galois multiplicity balance; it does
not copy the quadratic assertion that Galois squares preserve every
relative twist-trace type. A surviving linear system is retained as a
scoped necessary relaxation and does not supply a category.

